\subsection{A sharp asymptotic family of harmonic selected values}

\begin{theorem}[DP51: the two-column harmonic rate]\label{DP51}
Use DP46's arbitrary-unit singleton/complement pair over the unramified ring
$R=W(k)$, with $p$ odd and $h\ge3$.
Let $d=h-1$, $s=1/(p-1)$, and write $\beta$ for the primitive logarithm in additive coordinates.
Fix $1\le j\le d$ and $\lambda\in R^\times$, and put
\[
 N=p^{h-j},\qquad S=\frac{N-1}{p-1},\qquad
 w=\frac{sN}{N+1},\qquad c=\frac{s}{N+1},
 \qquad F(x,y)=\exp\bigl(\lambda\beta(x,e_jy)\bigr).
\]
For $a,b\ge0$ define the integral polynomial
\[
 P_{a,b}=\sum_{i=0}^{\min(a,b)}(-1)^i
 \binom ai\binom bi
 X_1^{a-i}Y_1^{b-i}X_2^iY_2^i
\]
and its actual primitive-kernel value
\[
 H_{a,b}=P_{a,b}(\partial)
       \bigl(F(x_1,y_1)F(x_2,y_2)\bigr)\big|_0.
\]
Then $P_{a,b}$ has fixed row degrees $(a,b)$ and is killed by
$\Delta=\partial_{X_1}\partial_{Y_1}+\partial_{X_2}\partial_{Y_2}$.
Moreover
\begin{equation}\label{eq:harmonicsharp}
 \vp(H_{a,b})\ge c(a+b)-\frac{s_p(a)+s_p(b)}{p-1},
 \qquad
 \liminf_{n\to\infty}\min_{a+b=n}\frac{\vp(H_{a,b})}{n}=c.
\end{equation}
Zero values have valuation $+\infty$; infinitely many nonzero values attain the limiting slope.
Any sequence attaining it necessarily satisfies
\[
 \frac{a}{a+b}\longrightarrow\frac1{N+1},\qquad
 \frac{b}{a+b}\longrightarrow\frac{N}{N+1}.
\]

For $j=1$, give the positive type one row, the negative type $d$ rows, and use only its first row.
Then these polynomials satisfy every algebraic condition in the recovered \path{HaO}: summed harmonicity, row invariance, fixed row degrees and integral coefficients.
The sharp slope on that family is $1/((p-1)(p^{h-1}+1))$.
For general $j$ the harmonic statement holds, but a full \path{HaO} interpretation requires that the chosen direction be a permitted first row; no reordering of a prescribed geometric datum is assumed.
The values are unchanged by balanced quadratic gauge factors lying in the formal ideal generated by the column-sum contractions that annihilate $P_{a,b}$.
\end{theorem}

\begin{proof}
Write $z(x,y)=\lambda\beta(x,e_jy)$.
We first recover the entire two-variable initial form from DP42 and DP46.
For every rational
\[
 t\in I=\left[\frac{S}{N+1},w\right],
\]
their contraction profile gives
\begin{equation}\label{eq:harmonicinitial}
 \nu_{(t,t)}(z)=(N+1)t-S,
\end{equation}
with a unique least monomial, a unit multiple of $xy^N/p^S$.
Indeed replace the scalar input $uT$ in those proofs by an independent variable $X$ and the surviving dual input $v_jT$ by $Y$.
The scalar inverse reduces to $X$.
The reduced dual inverse is zero before $j$ and then has successive coordinates equal to unit multiples of $Y,Y^p,\ldots,Y^{p^{d-j}}$.
The full primitive logarithm has the unique least group-coordinate term $S_{\rm group}Y_d^p/p$.
Composition therefore gives a unit multiple of $XY^N$, while all other terms remain strictly higher by the same integral-substitution reduction.
The common contraction powers obey $\mu=\gamma^{N+1}/p^S$.
This proves \eqref{eq:harmonicinitial}, including the coefficient valuation $-S$, with the arbitrary unit arrows and $\lambda$ included.
Invariant integral-frame corrections are strictly higher by DP46.

Choose a rational $t_0$ in the interior of $I$ below $w$.
Boundedness at $(t_0,t_0)$ gives, for every supported exponent,
\[
 \vp([x^ay^b]z)+w(a+b)
 \ge (N+1)t_0-S+(w-t_0)(a+b).
\]
These quantities tend to $+\infty$ with total degree.
At $(w,w)$ the unique least monomial $xy^N$ has weight $s$.
There is a small positive rational $\varepsilon$ such that both vectors
\begin{equation}\label{eq:harmonicperturb}
 (w-N\varepsilon,w+\varepsilon),\qquad
 (w+N\varepsilon,w-\varepsilon)
\end{equation}
have positive entries and give $z$ infimum $s$, with that same unique least monomial.
For rigor, first require $N\varepsilon<(w-t_0)/2$.
Every perturbed weight is then at least $t_0+(w-t_0)/2$, so all sufficiently large degrees have weight greater than $s+1$, uniformly over these perturbations.
Only finitely many other monomials remain; shrink $\varepsilon$ to preserve their strict inequalities.
The weight of $xy^N$ is unchanged because $-N\varepsilon+N\varepsilon=0$.
Thus no unjustified uniform-gap assumption about an infinite series is needed.

Since $\vp(m!)\le ms$, exponentiating a zero-constant series of coefficient infimum at least $s$ gives an integral series after contraction.
The inverse exponential is integral as well.
Therefore $F$ is integral with constant one at both vectors \eqref{eq:harmonicperturb} and at $(w,w)$.

Write $f_{a,b}=[x^ay^b]F$ and set
\[
 D(T)=\sum_{i\ge0}f_{i,i}T^i,\qquad
 C(x,y)=F(x,y)D(-xy)=\sum_{a,b\ge0}C_{a,b}x^ay^b.
\]
These formal series are well defined and $D(0)=1$.
The first vector in \eqref{eq:harmonicperturb} gives
\[
 \vp(f_{i,i})+i\bigl(2w-(N-1)\varepsilon\bigr)\ge0.
\]
Consequently $D(-T^2)$ is integral at the strictly smaller weight
\[
 w_D=w-\frac{N-1}{2}\varepsilon<w.
\]
Also $D(-xy)$ is integral at both vectors \eqref{eq:harmonicperturb}, whose coordinate sums are
$2w-(N-1)\varepsilon$ and $2w+(N-1)\varepsilon$.
Thus $C$ is integral at both vectors.

The first contraction of the $i$th summand of $P_{a,b}$ cancels the second contraction of its $(i+1)$st summand, because
\[
 \binom ai\binom bi(a-i)(b-i)
 =\binom a{i+1}\binom b{i+1}(i+1)^2.
\]
The end terms vanish when an exponent is zero.
This proves $\Delta P_{a,b}=0$.
Factorial coefficient extraction and
\[
 \binom ai\binom bi(a-i)!(b-i)!i!^2=a!b!
\]
give the exact generating identity
\begin{equation}\label{eq:harmonicextract}
 H_{a,b}=a!b!\sum_i(-1)^if_{a-i,b-i}f_{i,i}
         =a!b!C_{a,b}.
\end{equation}
Every harmonic correction term has been retained in this identity.

On the diagonal,
\[
 C(T,T)=F(T,T)D(-T^2).
\]
At rational weights sufficiently close to $w$ from below,
\eqref{eq:harmonicinitial} puts the logarithm of $F(T,T)$ in $[0,s)$, with initial form a unit times $T^{N+1}$.
Its derivative in the residue field is nonzero because $N+1$ is prime to $p$.
The factor $D(-T^2)$ is nonzero and bounded at every weight at least $w_D$.
DP41's bounded-multiplier lemma therefore makes $C(T,T)$ unbounded at all those weights below $w$.
Monotonicity gives unboundedness at every smaller nonnegative weight.
At and above $w$ both factors are integral with constant one.
Thus the diagonal has exact coefficient threshold $w$.

If the two-variable coefficients $C_{a,b}$ were bounded at $(t,t)$ for some $t<w$, their finite sums at each total degree would make the diagonal bounded there, a contradiction.
They are integral at $(w,w)$, so
\[
 \liminf_{n\to\infty}\min_{a+b=n}\frac{\vp(C_{a,b})}{n}=-w.
\]
A strictly larger lower limit would give boundedness at a sufficiently close smaller rational weight, which proves the equality explicitly.
Legendre's identity
\[
 \vp(a!b!)=s(a+b)-s\bigl(s_p(a)+s_p(b)\bigr)
\]
and the uniform $O(\log(n+1))$ digit-sum bound for $a+b=n$ prove
\eqref{eq:harmonicsharp} from \eqref{eq:harmonicextract}.

Both inequalities from \eqref{eq:harmonicperturb} further give
\[
 \vp(C_{a,b})\ge-w(a+b)+\varepsilon|Na-b|.
\]
Hence any sequence attaining the slope $c$ has
$|Na-b|/(a+b)\to0$, proving the asserted degree ray.

For the source conditions, use \path{RowDeg.lean}: \path{lap} lines 21--25,
\path{harmK} 27--30, \path{RowInv} 37--38, \path{RowDeg} 40--41 and
\path{HaO} 46--48; \path{IsThRule.lean}, \path{harmA} 86--87, specializes them.
For $j=1$ the first negative-row contraction is exactly $\Delta$, and all later-row contractions vanish.
Adding an earlier negative row to a later one leaves $P_{a,b}$ unchanged because it involves no later row.
The sole positive row has no row operation.
Each monomial has row degrees $a,b,0,\ldots,0$, and all coefficients are integers.
This proves precisely the stated algebraic \path{HaO} membership.

Finally DP39's Fischer identity makes the pairing with $P_{a,b}$ vanish on the formal ideal generated by the column-sum quadratics whose differential operators annihilate it.
If a zero-constant $Q$ lies in that ideal, so does $\exp(Q)-1$.
Multiplication by this balanced gauge factor therefore does not change the selected value.
This proves the gauge assertion with its exact scope.
\end{proof}

\begin{corollary}[DP51: the sharp universal polynomial bound]
Put $N_*=p^{h-1}$, $w_*=sN_* /(N_*+1)$ and $c_*=s/(N_*+1)$.
For $\ell$ independent copies, integral formal functions $U_b$, and units $\lambda_b$, set
\[
 \Theta(\mathbf Z)=\prod_{b=1}^{\ell}
 U_b(E_G(x_b),E_H(\mathbf y_b))
 \exp\bigl(\lambda_b\beta(x_b,\mathbf y_b)\bigr),
\]
with $h$ additive variables in each copy.
For every nonzero integral homogeneous differential polynomial $P$ of degree $n$,
\begin{equation}\label{eq:harmonicuniversal}
 \vp\bigl(P(\partial)\Theta(0)\bigr)
 \ge c_*n-s\max_{\alpha\in\operatorname{supp}P}
                  \sum_{i=1}^{\ell h}s_p(\alpha_i).
\end{equation}
For $n\ge1$ this implies the bound
$c_*n-\ell h(1+\floor{\log_p n})$.
For fixed $h$ and $\ell\ge1$, the slope $c_*$ is optimal uniformly over these integral polynomials and prefactors.
For $\ell\ge2$ it is already optimal on the source-shaped \path{HaO} submodule with one positive row and $d$ negative rows.
If the number of columns varies with $n$, the exact digit-sum error must be kept; no sublinear error uniform in a growing $\ell$ is asserted.
\end{corollary}

\begin{proof}
Use the profile of DP42/DP46 with $j=1$ at $w=w_*$.
Its scalar and first dual weights are $w_*$, and for $i>1$ its dual weights are
\[
 \omega_i=s-p^{i-1}(s-w_*)<w_*.
\]
All weights are positive and exceed the bare contraction weights.
The scaled inverse logarithms are integral, and the full group-coordinate primitive logarithm has coefficient infimum $s$ at this profile.
Give every independent additive variable the common weight $w_*$.
Its input into each contracted variable has nonnegative extra weight: zero on the scalar and first dual coordinates, and $w_*-\omega_i$ on later dual coordinates.
The inverse substitutions therefore remain integral and cannot decrease the primitive logarithm's infimum below $s$.
Its exponential with unit parameter is integral on the full polydisc.
Every integral prefactor remains integral after these substitutions as well.
Multiplying the copies gives
\[
 \vp([\mathbf Z^\alpha]\Theta)+w_*|\alpha|\ge0.
\]
The invariant derivative of multiorder $\alpha$ is $\alpha!$ times this additive coefficient.
Legendre's formula gives
$\vp(\alpha!)=s|\alpha|-s\sum_i s_p(\alpha_i)$.
The coefficients of $P$ are integral and all supported multiindices have degree $n$, so the nonarchimedean triangle inequality proves \eqref{eq:harmonicuniversal}.
For $0\le\alpha_i\le n$ and $n\ge1$ one has
$s_p(\alpha_i)\le(p-1)(1+\floor{\log_p n})$, giving the stated simpler error.

For the full polynomial module with one copy, take prefactor one and $P=(X+Y_1)^n$.
Its value is the ordinary derivative along $x=t,y_1=t$, whose exact slope is $c_*$ by DP42/DP46.
Extra copies with prefactor one do not affect that value.
For sharpness on the $\ell\ge2$ \path{HaO} submodule take two equal kernels with prefactors one and the polynomials $P_{a,b}$ above with $j=1$.
All unused variables are set to zero in their selected value; extra copies with prefactor one contribute their constant value one.
Their values are precisely $H_{a,b}$ and attain $c_*$ by \eqref{eq:harmonicsharp}.
Their full algebraic \path{HaO} membership was proved above.
A larger uniform slope with sublinear error is therefore impossible on either module.
This does not assert that every prescribed geometric prefactor or weight sequence attains that minimum.
\end{proof}

The full-series bound uses the stated invariant product normalization.
A balanced quadratic gauge factor killed by the harmonic contractions leaves the selected value unchanged, as proved above; an arbitrary rational quadratic before selection is not covered by that assertion.

DP51 closes an asymptotic local shortcut left open by a single-jet example:
the stated harmonic conditions alone do not give a larger uniform slope on this family.
It does not identify the pure kernel germ with a global theta section, force the prescribed global row degrees to approach the displayed ray, or decide the complete gain/cost inequality for a shorter chain.
Those questions still require the actual section, selected weights, path, period and global links.

